% === Conclusion ===
\section{Conclusion}
\label{sec:conclusion}

Incident triage is, at its core, an exercise in analogy: the fastest route from a
new alert to a resolution is to recognize which past incident it resembles. We
built \sys, an agentic triage system that operationalizes this view---retrieving
similar resolved incidents from a memory of past tickets by fusing dense, sparse,
and graph signals, adapting which reasoning steps it runs to the evidence each
window carries, and suppressing duplicate pages across an incident's lifetime.
Across a realism spectrum from fully instrumented synthetic services to a large
corpus of real Apache Jira incidents, hybrid retrieval locates the right prior
incident in its top five $97\%$ of the time on real data---outperforming strong
dense, lexical, graph, and LLM baselines---while the capability-gated controller
matches that quality at up to $78\%$ lower cost and degrades gracefully to
text-only operation.

Our broadest finding is a caution about evaluation. The ticket-worthiness
classification that dominates synthetic operations benchmarks is trivially solved
on injected faults yet unlearnable on real ones, so it measures the data rather
than the method; the difficulty that actually transfers lies in retrieval. We
hope the datasets, harness, and baselines we release help the community study
triage where it is genuinely hard---on real incidents, under real evidence
constraints.

\section*{Data Availability}

All source code, the evaluation harness, the two collected telemetry datasets, and
the complete result set---including the prior-art and LLM baselines, ablations, and
significance tests---are released to support reproduction; an anonymized artifact
accompanies this submission for review. The World of Logs corpus is publicly
available from its original authors under its own terms.
